\subsection{The output of the application, panel by panel}
The phenotyping layers have two pictures: the deterministic plant of the model and the volumetric scan of a real plant by a camera that turns round it.

The layer draws the plant of the model for the chosen crop and day, with the height, the leaf area index, the roots and the cover.

\appfig{dph_0}{0.96}{Deterministic phenotyping: the plant of the model through the season.}

After the run the numbers of the plant (height, LAI, biomass) are given for the days of the season.

\appfig{dph_1}{0.96}{Run simulation: the phenotype for the settings of the user.}

The page is a guided tour of the six steps from the rig to the measured volume.

\appfig{vph_0}{0.96}{Volumetric phenotyping: the plant, the camera rig and the six steps of the scan.}

The arm turns, a view arrives at each position and the views are joined into one cloud of points.

\appfig{vph_1}{0.96}{Step 2: the camera turns round the plant; the views arrive and are joined.}

The simulation can be played step by step to see the point cloud grow.

\appfig{vph_2}{0.96}{Step 2 of 6 (the camera turns round the plant), with the moving simulation of the scan.}

The same rig scans the plant on seven dates of the season to follow its growth.

\appfig{vph_3}{0.96}{Step 3: growth from emergence to the end of the season.}

The volume of the point cloud (voxel counting) is compared with the model value, an error check of the method.

\appfig{vph_4}{0.96}{Step 5: volume and form measured on the joined cloud against the model for the same day.}

The measured volumes form the growth curve and the line is the deterministic model.

\appfig{vph_5}{0.96}{Step 7: through the season, a scan every ten days (24 of 24 done); the line is the model.}

The last panel states the inputs of the layer and the numbers that depend on it.

\appfig{vph_6}{0.96}{Step 8: what feeds this layer and what it feeds.}

